% Generated draft: requires mathematical and source review.
\documentclass[11pt]{article}
\usepackage[margin=25mm]{geometry}
\usepackage{fontspec}
\setmainfont{Latin Modern Roman}
\usepackage{amsmath,amssymb,mathtools,unicode-math}
\setmathfont{Latin Modern Math}
\usepackage{enumitem,needspace,xurl,hyperref}
\setlength{\parindent}{0pt}
\setlength{\parskip}{5pt}
\setlength{\emergencystretch}{3em}
\newcommand{\N}{\mathbb{N}}
\newcommand{\Z}{\mathbb{Z}}
\newcommand{\Q}{\mathbb{Q}}
\newcommand{\R}{\mathbb{R}}
\newcommand{\C}{\mathbb{C}}
\newcommand{\F}{\mathbb{F}}
\newcommand{\A}{\mathbb{A}}
\newcommand{\PP}{\mathbb{P}}
\newcommand{\E}{\mathbb{E}}
\newcommand{\eps}{\varepsilon}
\newcommand{\dd}{\,\mathrm{d}}
\begin{document}
\section*{UnsolvedMath ID: 30001162}
\textbf{Source problem code:} catalogue entry
\subsection{Definitions and mathematical statement}
For notation, let \(d\in\N\) and let \(\Omega\subset\R^d\) be a nonempty bounded open set. Write
\[
V=|\Omega|,\qquad S=\mathcal H^{d-1}(\partial\Omega),\qquad \omega_d=|\{x\in\R^d:|x|<1\}|.
\]
Here \(\mathcal H^{d-1}\) denotes Hausdorff measure. Boundary-measure expressions below require \(S<\infty\). This is a framework for explaining the question; the supplied record does not specify the dimension, admissible domain class, or boundary regularity required for the proposed improvement.

Let
\[
0<\lambda_1^D(\Omega)\le\lambda_2^D(\Omega)\le\cdots
\]
be the eigenvalues of the nonnegative Dirichlet Laplacian \(-\Delta\), counted with multiplicity. The Li--Yau lower bound for their partial sums is
\[
\sum_{j=1}^{k}\lambda_j^D(\Omega)
\ge
\frac{d}{d+2}\frac{4\pi^2}{(\omega_d V)^{2/d}}\,k^{1+2/d},
\qquad k\in\N,\quad k\ge1.
\]
The supplied problem asks whether this bound admits a strengthening by a nonnegative, nontrivial correction involving the boundary surface measure \(S\). Schematically, the requested form is
\[
\sum_{j=1}^{k}\lambda_j^D(\Omega)
\ge
\frac{d}{d+2}\frac{4\pi^2}{(\omega_d V)^{2/d}}\,k^{1+2/d}
+\mathcal B_D(\Omega,k),
\]
where \(\mathcal B_D\) is to involve \(S\). This display records the direction of the requested improvement; it does not specify a proposed formula or assert that such an inequality holds. The source record leaves unresolved the allowable additional geometric dependence, the size and strict positivity of the correction, its uniformity over domains, and whether it should hold for every \(k\) or only in a specified regime.

Where the Neumann Laplacian has discrete spectrum, write its eigenvalues, including all zero eigenvalues and counting multiplicity, as
\[
0=\mu_1^N(\Omega)\le\mu_2^N(\Omega)\le\cdots.
\]
The problem also requests analogous boundary-measure corrections for sums of Neumann eigenvalues and for Berezin-type bounds, \emph{where appropriate}. For reference, the usual order-one Dirichlet Riesz mean is
\[
R_D(z)=\sum_{j\ge1}(z-\lambda_j^D(\Omega))_+,
\qquad (t)_+=\max\{t,0\},\qquad z\ge0,
\]
and its standard Berezin upper bound is
\[
R_D(z)\le
\frac{2}{d+2}\frac{\omega_d}{(2\pi)^d}\,V z^{1+d/2}.
\]
A strengthening of this particular upper bound would subtract a nonnegative boundary correction. However, the supplied record does not identify the intended Berezin-type functional or its order, nor the precise Neumann baseline, inequality direction, or parameter ranges. Consequently, these extensions remain requests for a precise formulation, rather than specified universally quantified inequalities. No connectedness assumption, boundary regularity hypothesis, or particular correction term is supplied.
\subsection{Short English statement}
Can the Dirichlet Li--Yau lower bound for sums of Laplacian eigenvalues be improved by a term involving boundary surface area? The record also asks for analogous Neumann and Berezin-type improvements where appropriate. It does not specify the correction, admissible domains, or all parameter ranges, so the displayed target is schematic and requires human mathematical review.
\subsection{Sources}
\begin{itemize}
\item Supplied local UnsolvedMath registry record, permanent ID \(30001162\), title \emph{Boundary Corrections to Li--Yau Eigenvalue Bounds}, source problem code \texttt{catalogue entry}, in the Oberwolfach Reports Open Problems collection. The record gives proposed year 2009 and names no individual proposer.
\item Supplied report citation: \emph{Low Eigenvalues of Laplace and Schrödinger Operators} (2009), Oberwolfach Reports, pp.\ 355--428, DOI \texttt{10.4171/owr/2009/06}. \url{https://doi.org/10.4171/owr/2009/06}. The supplied background describes these reports as records of mathematical workshops and their open-problem discussions; no precise problem page or verbatim report excerpt is supplied here.
\item Supplied literature source: Filonov--Levitin--Polterovich--Sher, described in the registry as ``Pólya for balls.'' \url{https://doi.org/10.1007/s00222-023-01198-1}.
\item Literature assessment dated exactly \texttt{2026-08-21}: the supplied curation labels the problem open and states that the Neumann Pólya conjecture remains open for general Euclidean domains, and that the broader request for a universal Li--Yau correction involving boundary measure remains unresolved despite partial bounds. This is supplied dated curation, not an independently verified literature assessment; its status, cited claims, and applicability to the intended formulation require independent verification.
\item The registry labels the statement \texttt{corrected\_verified} and reports that its mathematical content was checked against the official Oberwolfach report, with notation, context, and question form normalized without changing the source-level claim. This is supplied verification metadata, not a verification performed for this draft.
\end{itemize}
\end{document}
